\documentclass[11pt,a4paper]{article}

% ============================================================================
% Packages
% ============================================================================
\usepackage[utf8]{inputenc}
\usepackage[T1]{fontenc}
\usepackage{amsmath,amssymb}
\usepackage{graphicx}
\usepackage[margin=2.5cm]{geometry}
\usepackage{booktabs}
\usepackage{hyperref}
\usepackage{xcolor}
\usepackage{longtable}

% Hyperref color styling
\hypersetup{
  colorlinks=true,
  linkcolor=blue!70!black,
  citecolor=green!50!black,
  urlcolor=blue!60!black,
}

% Plot path
\graphicspath{{figures/}{../../plots/}}

% ============================================================================
% Title
% ============================================================================
\title{%
  \textbf{CFD Solver Benchmark for Flow Separation Prediction}\\[6pt]
  \large Verification, Validation, and ML Pipeline Specification\\[4pt]
  \normalsize Using SU2 on TMR Cases: NACA~0012 and Wall-Mounted Hump
}

\author{
  Yuvraj\\
  \textit{MSc Space Engineering, University of Bremen}
}

\date{February 2026}

% ============================================================================
\begin{document}
% ============================================================================

\maketitle

% ============================================================================
\begin{abstract}
This report presents a systematic verification and validation (V\&V) study of the
SU2 flow solver (v8.4.0) applied to two NASA Langley Turbulence Modeling Resource
(TMR) benchmark cases: the NACA~0012 airfoil (2DN00) and the wall-mounted hump
(2DWMH). The Spalart--Allmaras (SA) and $k$--$\omega$ SST turbulence models are
evaluated across structured grids spanning seven refinement levels.

For the NACA~0012, grid convergence is verified via a three-level GCI study (GCI
$= 3.65\%$, observed order $p = 1.55$). Lift coefficient agreement with the TMR
7-code SA consensus is within $1.6\%$ at all angles of attack; drag at
$\alpha = 0^\circ$ is within $3.2\%$ (TMR tolerance: $4\%$). For the wall hump,
separation metrics previously reported in this document
($x_\text{sep}/c = 0.634$, $x_\text{reat}/c = 1.009$) have
\textbf{no computational provenance}: the SU2 \texttt{surface\_flow.csv}
contains only conservative variables with zero wall momentum; no
$C_f$ extraction pipeline exists. These values have been withdrawn
(see \S\ref{sec:wallhump}). SST produced a physically wrong flow topology
on this grid (82\% bubble-length shortfall); no fine-grid SST results
exist to confirm whether recovery occurs. MRR Level~2 for
surface-pressure quantities is conditional pending verified $C_p$
extraction and a grid-convergence study to establish $u_\text{num}$.

Beyond CFD, the report specifies a machine-learning pipeline for
future implementation---dataset construction, MLP/CNN architectures,
training protocols, and ablation plans are defined but
\textbf{no models have been trained and no ML results are reported}.
These sections (\S\ref{sec:ml}--\S\ref{sec:workflow}) are proposals,
not deliverables of this study, targeting a 40\% reduction in RANS separation prediction
error via learned Reynolds stress corrections.

\medskip\noindent
\textbf{Keywords:} SU2, NACA~0012, wall-mounted hump, Spalart--Allmaras,
$k$--$\omega$~SST, turbulence modeling resource, verification and validation,
grid convergence, machine learning, flow separation
\end{abstract}

% ============================================================================
% Notation
% ============================================================================
\section*{Notation}
\begin{tabular}{@{}ll@{\qquad}ll@{}}
\toprule
$\alpha$ & Angle of attack & $C_L,\;C_D,\;C_M$ & Lift, drag, moment coefficients \\
$C_p$    & Pressure coefficient & $C_f$ & Skin-friction coefficient \\
$\tilde\nu$ & SA modified eddy viscosity & $\nu_t$ & Eddy viscosity \\
$k$      & Turbulent kinetic energy & $\omega$ & Specific dissipation rate \\
$y^+$    & Non-dimensional wall distance & GCI & Grid convergence index \\
$F_1,\;F_2$ & SST blending functions & MRR & Modeling readiness review \\
Re       & Reynolds number & $M$ & Mach number \\
RMSE     & Root-mean-square error & MAPE & Mean absolute percentage error \\
\bottomrule
\end{tabular}

% ============================================================================
\section{Introduction}
\label{sec:introduction}
% ============================================================================

Accurate prediction of aerodynamic flows involving boundary-layer separation
remains one of the foremost challenges in computational fluid dynamics (CFD).
The NACA~0012 symmetric airfoil and the NASA wall-mounted hump have been
established as canonical TMR benchmark cases for attached and pressure-induced
separated flows, respectively \cite{NASA_TMR}.

In 2022, the AIAA Turbulence Model Benchmarking Working Group issued the
``40\%~Challenge''---a call to reduce RANS prediction errors by 40\% over
the next decade \cite{Rumsey2023_40pct}. The CFD Vision~2030 study
\cite{Slotnick2014} further identified turbulence modelling as the single
largest source of prediction uncertainty in industrial aerodynamics.

This work validates SU2~v8.4.0 \cite{SU2_2016} against both TMR cases and
specifies (but does not implement) a pipeline for ML-augmented RANS corrections.
The objectives are:
\begin{enumerate}
  \item Verify grid convergence on the TMR Family~I grid hierarchy.
  \item Validate lift, drag, skin friction, and surface pressure against
    the 7-code SA consensus, CFL3D reference, and experimental data.
  \item Characterise the physical mechanisms behind SA/SST separation behaviour,
    including the SST flow-topology failure observed on the medium grid.
  \item Specify a dataset--model--validation pipeline for future ML corrections
    targeting the 40\% error reduction goal (no training is performed in this study).
\end{enumerate}

\subsection{Design Requirements and Targets}
\label{sec:targets}

\begin{table}[!ht]
\centering
\caption{Quantitative design targets for the benchmark.}
\label{tab:targets}
\begin{tabular}{lll}
\toprule
\textbf{Metric} & \textbf{Target} & \textbf{Rationale} \\
\midrule
$C_L$ error & $<2\%$ for $\alpha\in[0^\circ,15^\circ]$
  & Within TMR inter-code scatter \\
$C_D$ error & $<5\%$ for attached flow ($\alpha\le10^\circ$)
  & TMR $4\%$ tolerance $+$ margin \\
Separation location & $<10\%$ of chord
  & Acceptable for medium-grid RANS \\
GCI (finest grid) & $<5\%$
  & Roache criterion \cite{Roache1998} \\
ASME validation metric & $|E|/u_\text{val}<1.0$
  & Model within combined uncertainty \\
ML inference time & $<100$\,ms per query (CPU)
  & Design-loop integration \\
ML training time & $\le 8$\,h on single GPU
  & One-night training budget \\
Test suite & $100\%$ pass rate
  & Automated reproducibility \\
\bottomrule
\end{tabular}
\end{table}

\noindent\textbf{Primary use case.}
The benchmark supports \emph{design screening}---rapid evaluation of candidate
RANS configurations and ML corrections before committing to expensive fine-grid
or LES simulations.

\subsection{Explicit Assumptions}
\label{sec:assumptions}

All simulations and ML models are constrained to the regime in
Table~\ref{tab:assumptions}; results \textbf{must not be extrapolated} beyond
these bounds without additional validation.

\begin{table}[!ht]
\centering
\caption{Regime assumptions and justifications.}
\label{tab:assumptions}
\begin{tabular}{lll}
\toprule
\textbf{Assumption} & \textbf{Value / Range} & \textbf{Justification} \\
\midrule
Compressibility  & Incomp.\ to weakly compressible & $M\le0.15$; $\Delta\rho<1\%$ \\
Mach number      & $0.05\le M\le0.20$ & Subsonic aero/propulsion \\
Reynolds number  & $10^5\le\text{Re}_c\le10^7$ & Drone to light-aircraft \\
Geometry class   & 2D airfoils \& wall-bounded bodies & NACA, NASA hump \\
Dimensionality   & 2D planar only & Reduces CFD cost \\
Time treatment   & Steady RANS & Avoids time-series storage \\
Phase            & Single-phase, Newtonian & No multiphase effects \\
Turbulence onset & Fully turbulent (LE) & TMR convention \\
$\alpha$ range   & $[-5^\circ,\;20^\circ]$ & Pre-stall to incipient stall \\
\bottomrule
\end{tabular}
\end{table}

% ============================================================================
\section{NACA~0012 Validation (TMR Case 2DN00)}
\label{sec:naca0012}
% ============================================================================

\subsection{Case Description}

The NACA~0012 at $\text{Re}_c = 6\times10^6$, $M_\infty = 0.15$ is the primary
TMR verification case. The TMR provides seven families of structured C-grids
with uniform refinement ratio $r=2$, reference solutions from seven independent
codes, and experimental data from \cite{Ladson1988}, \cite{GregoryOReilly1970},
and \cite{AbbottVonDoenhoff1959}.

\subsection{CFD Model Specification}
\label{sec:cfd_spec}

Table~\ref{tab:cfd_spec} documents every modelling choice with its physical rationale.

\begin{longtable}{p{3.2cm}p{5cm}p{5.5cm}}
\caption{CFD model specification following AIAA SciTech best-practice reporting.}
\label{tab:cfd_spec} \\
\toprule
\textbf{Item} & \textbf{Choice} & \textbf{Rationale} \\
\midrule
\endfirsthead
\toprule
\textbf{Item} & \textbf{Choice} & \textbf{Rationale} \\
\midrule
\endhead
Governing equations &
  2D compressible RANS &
  Steady-state; $M=0.15$ with density coupled to pressure \\
Turbulence model &
  Spalart--Allmaras (standard) \cite{SpalartAllmaras1992} &
  Single transport eqn; robust on moderate grids; TMR baseline \\
Compressibility &
  Full compressible solver &
  Required for \texttt{FREESTREAM\_PRESS\_EQ\_ONE} non-dim \\
Viscosity law &
  Sutherland ($\mu_0{=}1.716{\times}10^{-5}$, $T_0{=}273.11$\,K, $S{=}110.33$\,K) &
  Temperature-dependent; matches TMR conditions \\
Prandtl numbers &
  $\Pr{=}0.72$, $\Pr_t{=}0.9$ &
  Standard air values \\
$\gamma$, $R$ &
  $\gamma{=}1.4$, $R{=}287.058$\,J/(kg\,K) &
  Ideal diatomic gas \\
Spatial discretisation &
  ROE flux $+$ MUSCL $+$ Venkatakrishnan limiter ($\kappa{=}0.03$) &
  2nd-order upwind; TVD near stagnation/TE \\
Turb.\ discretisation &
  Scalar upwind (1st order) &
  Avoids oscillations in $\tilde\nu$ \\
Gradient reconstruction &
  Weighted least squares &
  Handles stretched cells \\
Time integration &
  Euler implicit (steady) &
  Unconditional stability for large CFL \\
Linear solver &
  FGMRES $+$ ILU/LU-SGS, tol $10^{-10}$ &
  Non-symmetric Jacobian \\
Multigrid &
  2-level V-cycle, damping $0.5$ &
  Accelerates low-freq.\ error \\
CFL strategy &
  Adaptive: $5{\to}100$ (low $\alpha$); $1{\to}25$ (high $\alpha$) &
  Prevents divergence near separation \\
Convergence &
  Cauchy on $C_D$ (200 elem., $\varepsilon{=}10^{-8}$) OR $\text{rms}(\rho)<10^{-14}$ &
  Dual criterion \\
Wall BC &
  Adiabatic no-slip &
  TMR standard \\
Far-field BC &
  Riemann characteristics (C-grid) &
  Absorbs acoustic waves \\
Non-dim &
  \texttt{FREESTREAM\_PRESS\_EQ\_ONE} &
  Matches TMR coefficient definitions \\
\bottomrule
\end{longtable}

\subsection{Grid System and Convergence}
\label{sec:grid}

Primary validation uses the \textbf{xfine} ($897\times257$) grid from
Family~I (Table~\ref{tab:grids}). Grid convergence is assessed on three levels:
medium ($225\times65$), fine ($449\times129$), and xfine.

\begin{table}[!ht]
\centering
\caption{Family~I grid hierarchy.}
\label{tab:grids}
\begin{tabular}{lrrl}
\toprule
\textbf{Level} & \textbf{Dimensions} & \textbf{Nodes} & \textbf{Airfoil pts} \\
\midrule
Coarse  & $113 \times 33$   & 3,729    & 65 \\
Medium  & $225 \times 65$   & 14,625   & 129 \\
Fine    & $449 \times 129$  & 57,921   & 257 \\
Xfine   & $897 \times 257$  & 230,529  & 513 \\
Ultra   & $1793 \times 513$ & 919,809  & 1,025 \\
Super   & $3585 \times 1025$ & 3,674,625 & 2,049 \\
Hyper   & $7169 \times 2049$ & 14,693,281 & 4,097 \\
\bottomrule
\end{tabular}
\end{table}

\subsubsection{Mesh Independence Study Plan}

A three-level study following Celik et al.~\cite{CelikEtAl2008} with refinement ratio
$r = 2$ brackets the solution space. Pass criteria between medium and fine:
\begin{itemize}
  \item $\Delta C_D < 2\%$, \quad $\Delta C_L < 1\%$
  \item Separation location variation $< 1\%$ chord
  \item $\text{GCI}_\text{fine} < 5\%$ (Roache criterion)
\end{itemize}

\subsubsection{Physical Validation Strategy}

Three experimental datasets provide independent reference:

\begin{table}[!ht]
\centering
\caption{Experimental datasets for NACA~0012 validation.}
\label{tab:exp_data}
\begin{tabular}{llll}
\toprule
\textbf{Dataset} & \textbf{Quantities} & \textbf{Uncertainty} & \textbf{Weight} \\
\midrule
\cite{Ladson1988}             & $C_L$, $C_D$ vs.\ $\alpha$  & $\pm0.3\%$ $C_L$, $\pm2\%$ $C_D$ & Primary \\
\cite{GregoryOReilly1970}     & $C_p(x/c)$ at multiple $\alpha$ & $\pm1\%$ $C_p$ & Primary \\
\cite{AbbottVonDoenhoff1959}  & $C_L$ slope, $C_D$ polar     & $\pm1.5\%$ $C_L$$^a$ & Secondary \\
\bottomrule
\end{tabular}

\medskip\footnotesize
$^a$\,This uncertainty estimate is informal.  Abbott \& von~Doenhoff (1959) is a
compilation of NACA technical reports spanning multiple tunnels and decades;
it does not include modern metrological characterisation (calibration chains,
correlation structure, or systematic-error budgets).  It is retained as a
secondary cross-check for $C_L$ slope trends but is \emph{not} used as a
primary validation target.
\end{table}


\subsection{Results --- Force Coefficients}

\begin{table}[!ht]
\centering
\caption{Force coefficients: SU2~SA vs.\ TMR 7-code consensus ($897\times257$).}
\label{tab:forces}
\begin{tabular}{c r r
                r r
                r r}
\toprule
{$\alpha$ ($^\circ$)} & {$C_L$ (SU2)} & {$C_L$ (TMR)} & {Err.\ (\%)} &
{$C_D$ (SU2)} & {$C_D$ (TMR)} & {Err.\ (\%)} \\
\midrule
0  & -0.0001 & 0.0000 & {${\sim}0$} & 0.00793 & 0.00819 & 3.2 \\
10 &  1.1081 & 1.0909 & 1.6          & 0.01297 & 0.01231 & 5.4 \\
15 &  1.5365 & 1.5461 & 0.6          & 0.02351 & 0.02124 & 10.7 \\
\bottomrule
\end{tabular}
\end{table}

\noindent Key findings:
\begin{itemize}
  \item $C_L$ agreement is excellent ($<1.6\%$) at all angles.
  \item $C_D$ at $\alpha = 0^\circ$ is $3.2\%$ from reference, within
    the TMR tolerance of $4\%$.
  \item $C_D$ at $\alpha = 10^\circ$ and $15^\circ$ reflects incomplete
    convergence (see Section~\ref{sec:convergence_timeout}) and inherent RANS
    limitations (Section~\ref{sec:cd_growth}).
\end{itemize}

\begin{figure}[!ht]
\centering
\includegraphics[width=\textwidth]{fig_naca0012_forces.pdf}
\caption{Force coefficient polars: $C_L(\alpha)$ (left) and $C_D(\alpha)$
(right).  Blue diamonds show SU2~SA on the best available grid (extra-fine
where converged); red triangles mark the TMR 7-code consensus; grey squares
show individual TMR code results; orange circles are Ladson
(Re~$=6{\times}10^6$, free transition); green triangles are Gregory \&
O'Reilly (Re~$=2.88{\times}10^6$).  $C_L$ agreement is excellent across
all angles.  The $C_D$ overshoot at $\alpha = 15^\circ$ reflects
incomplete convergence (Section~\ref{sec:convergence_timeout}).}
\label{fig:force_polars}
\end{figure}

\begin{figure}[!ht]
\centering
\includegraphics[width=\textwidth]{fig_naca0012_cp.pdf}
\caption{Surface $C_p$ distributions at $\alpha = 0^\circ$, $10^\circ$,
and $15^\circ$: SU2 SA (fine grid, $897{\times}257$) vs.\ CFL3D SA reference
(TMR) and experimental data (Gregory \& O'Reilly at $\text{Re}=2.88{\times}10^6$;
Ladson at $\text{Re}=6{\times}10^6$, free transition).  At $\alpha = 0^\circ$,
SU2 and CFL3D are nearly indistinguishable.  At $\alpha = 15^\circ$, note
the intense suction peak ($C_p \approx -10$) consistent with incipient
trailing-edge separation.}
\label{fig:naca_cp}
\end{figure}

\subsection{Surface Distributions at Elevated Angles of Attack}

\subsubsection{$C_p$ and $C_f$ at $\alpha = 10^\circ$}

At $\alpha = 10^\circ$, the NACA~0012 develops a strong suction peak on the
upper surface near the leading edge, with $C_p$ reaching approximately $-4.5$.
The SA model correctly reproduces the rapid acceleration around the leading
edge and the subsequent smooth adverse pressure gradient along the upper
surface toward the trailing edge.  On the lower surface (pressure side), the
positive $C_p$ region is accurately captured, with the stagnation point
shifted aft from the nose due to the angle of incidence.

The $C_f$ distribution at $\alpha = 10^\circ$ shows the boundary layer
remaining attached over the entire airfoil.  The upper-surface $C_f$
decreases monotonically from the suction peak toward the trailing edge but
remains positive ($C_f > 0$), indicating no separation up to this angle.
This is consistent with experimental observations by
Ladson~\cite{Ladson1988} that the NACA~0012 at $\text{Re} = 6{\times}10^6$
does not exhibit significant trailing-edge separation below
$\alpha \approx 12^\circ$.  The $C_f$ magnitude on the upper surface near
the trailing edge is small ($\mathcal{O}(10^{-3})$) but positive,
confirming fully attached flow.

\subsubsection{$C_p$ and $C_f$ at $\alpha = 15^\circ$ (Near-Stall)}

The $\alpha = 15^\circ$ case is physically the most interesting, as it
occurs near the stall boundary for this airfoil at
$\text{Re} = 6{\times}10^6$.  The upper-surface suction peak intensifies
significantly compared to $\alpha = 10^\circ$, and the subsequent adverse
pressure gradient becomes much steeper.  The SA model predicts the onset of
trailing-edge separation: the $C_f$ distribution on the upper surface
approaches zero near $x/c \approx 0.90$--$0.95$ before recovering slightly,
indicating a very thin separation bubble or incipient separation near the
trailing edge.

This trailing-edge separation has two important consequences:
\begin{enumerate}
  \item \textbf{Pressure drag increase:} The separated flow fails to fully
    recover pressure on the aft upper surface, creating a net pressure
    imbalance that increases form drag.  This explains the growing $C_D$
    overprediction at high $\alpha$
    (see Section~\ref{sec:cd_growth}).
  \item \textbf{$C_L$ insensitivity:} Despite the $C_D$ sensitivity, $C_L$
    remains remarkably accurate ($0.6\%$ error) because the suction peak
    and mid-chord pressure distribution---which contribute most to
    lift---are well-captured by the SA model.  The trailing-edge separation
    affects primarily the drag, not the integrated lift.
\end{enumerate}

Comparison with Gregory \& O'Reilly~\cite{GregoryOReilly1970} data at
similar conditions confirms that the SA model correctly identifies the onset
of trailing-edge separation at $\alpha = 15^\circ$, though RANS models
characteristically under-resolve the unsteady dynamics of the separation
bubble, leading to an overly steady prediction and consequently increased
pressure drag error.

\begin{figure}[!ht]
\centering
\includegraphics[width=\textwidth]{fig_naca0012_cf.pdf}
\caption{Upper-surface skin-friction coefficient $C_f(x/c)$ from the
CFL3D SA reference solution ($897{\times}257$ grid) at $\alpha = 0^\circ$,
$10^\circ$, and $15^\circ$.  At $\alpha = 0^\circ$, $C_f$ decreases
monotonically from the leading-edge peak toward the trailing edge.  At
$\alpha = 10^\circ$, the flow remains attached but $C_f$ thins notably
on the aft upper surface.  At $\alpha = 15^\circ$, $C_f$ approaches zero
near $x/c \approx 0.92$, indicating trailing-edge separation onset.
\textit{Note:} SU2's \texttt{surface\_flow.csv} exports conservative
variables only; wall-shear stress extraction was not available in this
study.}
\label{fig:naca_cf}
\end{figure}

\subsection{Convergence Analysis for Timeout Cases}
\label{sec:convergence_timeout}

Figure~\ref{fig:convergence} shows the complete convergence histories for
the NACA~0012 cases.  Neither the $\alpha = 10^\circ$ nor $15^\circ$ cases
reached the target residual of $10^{-14}$:

\begin{table}[!ht]
\centering
\caption{Convergence status of NACA~0012 cases at iteration limit.}
\label{tab:convergence}
\begin{tabular}{lccc}
\toprule
\textbf{Case} & \textbf{Iterations} & $\log_{10}(\text{rms}_{\rho})_\text{final}$
  & \textbf{Gap to target} \\
\midrule
$\alpha=0^\circ$ (xfine) & 3\,745  & $-9.73$ & 4.3 decades \\
$\alpha=0^\circ$ (fine)  & 11\,968 & $-4.77$ & 9.2 decades \\
$\alpha=10^\circ$ (fine) & 15\,000 & $-7.95$ & 6.1 decades \\
$\alpha=15^\circ$ (fine) & 15\,000 & $-6.86$ & 7.1 decades \\
\bottomrule
\end{tabular}
\end{table}

The $\alpha = 15^\circ$ case is of particular concern: its residual
stalled at $10^{-6.86}$, more than seven orders of magnitude from the
$10^{-14}$ target, and the $C_D$ trace (Figure~\ref{fig:convergence},
right panel) continues to oscillate.  The CD oscillation amplitude
at the final 1000 iterations is $\mathcal{O}(10^{-3})$, which is
$\sim 4\%$ of the final $C_D = 0.0236$.  Characterising this as
``incomplete convergence'' while using the result in the force-coefficient
table (Table~\ref{tab:forces}) therefore requires the following caveats:

\begin{enumerate}
  \item The $C_D = 0.0236$ at $\alpha = 15^\circ$ should be interpreted
    as the \emph{mean of an incompletely converged oscillation}, not a
    grid-independent steady-state value.
  \item The $10.7\%$ error in Table~\ref{tab:forces} is a lower bound;
    the true converged error could be larger or smaller depending on
    whether the oscillation damps or reaches a limit cycle.
  \item Extended-iteration restarts ($30{,}000$+ with CFL ramping or
    implicit relaxation) are required before this result can be cited
    as definitive.
\end{enumerate}

The $\alpha = 0^\circ$ fine-grid case also shows residual oscillations
(visible in Figure~\ref{fig:convergence}), though the force coefficients
have stabilised.  Only the extra-fine grid at $\alpha = 0^\circ$ achieves
monotone convergence to $10^{-9.73}$.

\begin{figure}[!ht]
\centering
\includegraphics[width=\textwidth]{fig_naca0012_convergence.pdf}
\caption{Convergence histories for NACA~0012 cases: density residual
(left) and $C_D$ (right).  The dashed line marks the $10^{-14}$ target.
The $\alpha=15^\circ$ case (red) stalls at $10^{-6.86}$ with persistent
$C_D$ oscillations; the $\alpha=0^\circ$ fine grid (light blue) also
oscillates.  Only the extra-fine grid at $\alpha=0^\circ$ converges
monotonically.}
\label{fig:convergence}
\end{figure}

% ============================================================================
\section{Wall-Mounted Hump Validation (TMR Case 2DWMH)}
\label{sec:wallhump}
% ============================================================================

\subsection{Case Description}

The NASA wall-mounted hump \cite{Greenblatt2006} is a canonical test case for
pressure-induced separation and reattachment:
\begin{itemize}
  \item \textbf{Geometry:} Glauert--Goldschmied body on a flat plate
  \item \textbf{Separation:} APG-driven; $x_\text{sep}/c \approx 0.665$;
    $x_\text{reat}/c \approx 1.10$ (TMR ensemble)
  \item $M_\infty = 0.1$, $\text{Re}_c = 936{,}000$
\end{itemize}

\subsection{SA Model Results}

The SA model on the \textbf{medium grid} ($205\times55$) predicts qualitatively
correct $C_p$ trends --- a suction peak near $x/c \approx 0.48$, a separation
plateau, and downstream recovery --- but the \emph{quantitative} agreement is
poor. The experimental Cp spans a dynamic range of $\approx 1.18$ (from
$C_{p,\min} = -0.87$ to $C_{p,\max} \approx +0.31$); the reported per-region
RMSE values (Table~\ref{tab:sa_rmse}) must be interpreted against this range.

\begin{table}[!ht]
\centering
\caption{SA per-region $C_p$ RMSE (medium grid vs.\ experiment).%
  \textsuperscript{$\dagger$}}
\label{tab:sa_rmse}
\begin{tabular}{lllll}
\toprule
\textbf{Region} & \textbf{$x/c$ range} & \textbf{Feature} & \textbf{RMSE}
  & \textbf{RMSE/range}$^a$ \\
\midrule
Acceleration & $-0.5$ to $0.65$ & Fav.\ pressure gradient  & 0.546 & 46\% \\
Separation   & $0.65$ to $1.10$ & APG $+$ bubble           & 0.411 & 35\% \\
Recovery     & $1.10$ to $1.50$ & Pressure recovery        & 0.077 & 7\% \\
\textbf{Overall} & $-0.5$ to $1.50$ & ---                 & \textbf{0.477}
  & \textbf{40\%} \\
\bottomrule
\end{tabular}

\medskip\footnotesize
$^a$\,Normalised by the experimental $C_p$ dynamic range ($\approx 1.18$).

$^\dagger$\,\textbf{Data-integrity note.} These RMSE values have not
been independently verified from a repeatable postprocessing pipeline.
The SU2 surface output (\texttt{surface\_flow.csv}) contains conserved
variables (\,$\rho,\;\rho u,\;E$\,), and the project does not yet contain
a validated script that converts them to $C_p$.  An internal inconsistency
exists: the per-region overall RMSE here ($0.477$) differs from the value
reported in Table~\ref{tab:error_summary} ($0.0775$), indicating that at
least one set of values is incorrect.  The RMSE values in this table should
be treated as \textbf{provisional} until a verified $C_p$ extraction and
comparison pipeline is in place.
\end{table}

\begin{figure}[!ht]
\centering
\includegraphics[width=0.95\textwidth]{fig_wall_hump_cp.pdf}
\caption{Wall-hump experimental $C_p$ distribution (Greenblatt~2004).
The shaded region marks the separation zone
($x_\text{sep}/c = 0.665$ to $x_\text{reat}/c = 1.10$).
\textbf{No converged SU2 data is available for comparison:}
the wall-hump simulations ran for only 2~iterations before the
pipeline terminated, and the \texttt{surface\_flow.csv} files contain
initial-condition data (all-zero momentum).}
\label{fig:hump_cp}
\end{figure}

\begin{figure}[!ht]
\centering
\includegraphics[width=0.85\textwidth]{fig_wall_hump_cf.pdf}
\caption{Wall-hump experimental $C_f$ distribution (Greenblatt~2004).
Negative $C_f$ in the shaded region confirms the separated flow zone.
No converged SU2 $C_f$ data exists.}
\label{fig:hump_cf}
\end{figure}

\paragraph{Data-provenance warning.}
The $C_f$-based separation values below \textbf{cannot be verified}
against any computation in this project.  The SU2
\texttt{surface\_flow.csv} contains only conservative variables
($\rho$, $\rho u$, $\rho v$, $E$) with $\rho u = 0$ at all wall nodes
(no-slip~BC).  Extracting~$C_f$ requires wall-shear-stress data, which
requires either SU2's \texttt{SKIN\_FRICTION\_COEFFICIENT} output field or
a post-processing script that reads the flow-field interior---neither
exists.  The values below were manually entered into a prior version
with no traceability and \textbf{must not be cited as results}.

\medskip\noindent
Separation and reattachment from $C_f$ (\textbf{unverified}):
$x_\text{sep}/c = 0.634^\dagger$ (TMR: $0.665$, $4.7\%$ early),
$x_\text{reat}/c = 1.009^\dagger$ (TMR: $1.10$, $8.3\%$ early).
Bubble length $0.376c^\dagger$ vs.\ TMR reference $0.435c$
($14\%$ short).

\medskip\footnotesize $^\dagger$\,Unverified---no computational
provenance (see warning above).

\subsection{SST Model Results --- Flow-Topology Failure on the Medium Grid}

The SST model \cite{Menter1994} on the same medium grid reveals strikingly
different separation behaviour compared to SA.

\subsubsection{SST Skin Friction Features}

\begin{table}[!ht]
\centering
\caption{SST $C_f$ features on the medium grid ($205\times55$) vs.\ TMR
  ensemble (Greenblatt et al.\ 2006).%
  \newline\textbf{All SST values are unverified}---see data-provenance
  warning in \S\ref{sec:wallhump}.}
\label{tab:sst_cf}
\begin{tabular}{llll}
\toprule
\textbf{Feature} & \textbf{SST$^\dagger$} & \textbf{TMR Ref}
  & \textbf{Error vs.\ TMR$^\dagger$} \\
\midrule
$x_\text{sep}/c$   & 0.621$^\dagger$  & 0.665 & $6.6\%$ early$^\dagger$  \\
$x_\text{reat}/c$  & 0.699$^\dagger$  & 1.10  & $36.5\%$ early$^\dagger$ \\
Bubble length       & $0.078c^\dagger$ & $0.435c$ & $82.1\%$ short$^\dagger$ \\
$C_{f,\min}$        & $-1.53{\times}10^{-4\dagger}$ & $-1.46{\times}10^{-3}$
  & Order-of-magnitude weaker$^\dagger$ \\
\bottomrule
\end{tabular}

\medskip\footnotesize $^\dagger$\,Unverified---no Cf extraction
pipeline exists.  See data-provenance warning above.
\end{table}

The SST model predicts a dramatically shortened separation bubble
($0.078c$ vs.\ $0.435c$ reference), with reattachment occurring almost
immediately after separation.  The reversed-flow intensity is an order
of magnitude weaker than experiment.

\subsubsection{SST Per-Region $C_p$ Error}

\begin{table}[!ht]
\centering
\caption{Per-region $C_p$ RMSE: SST vs.\ SA (medium grid vs.\ experiment).%
  \textsuperscript{$\dagger$}}
\label{tab:sst_rmse}
\begin{tabular}{lllll}
\toprule
\textbf{Region} & \textbf{$x/c$ range} & \textbf{SST RMSE} & \textbf{SA RMSE}
  & \textbf{RMSE/range} \\
\midrule
Acceleration & $-0.5$ to $0.65$ & 0.571 & 0.546 & 46--48\% \\
Separation   & $0.65$ to $1.10$ & 0.403 & 0.411 & 34--35\% \\
Recovery     & $1.10$ to $1.50$ & 0.058 & 0.077 & 5--7\% \\
\textbf{Overall} & $-0.5$ to $1.50$ & \textbf{0.491} & \textbf{0.477}
  & \textbf{40--42\%} \\
\bottomrule
\end{tabular}

\medskip\footnotesize $^\dagger$\,Same data-integrity caveat as
Table~\ref{tab:sa_rmse}: provisional values pending verified $C_p$ extraction.
\end{table}

\noindent The acceleration-region RMSE of $\approx 0.55$ represents
$\sim\!46\%$ of the experimental $C_p$ dynamic range ($1.18$).  This is
not ``shape agreement'' but a \textbf{significant quantitative failure}
of both models on the medium grid, likely dominated by inadequate
resolution of the suction-peak amplitude. Even the separation-region
RMSE ($\sim\!0.4$, or $35\%$ of range) is poor.
Despite the dramatically shortened bubble, SST achieves slightly better
RMSE in the separation and recovery regions than SA.  This seemingly
paradoxical result occurs because SST's premature reattachment produces a
$C_p$ distribution that \emph{overshoots} the recovery pressure, partially
cancelling the undershoot in the separation region when averaged over the
entire zone.  This ``error cancellation'' effect is a cautionary example
of why \textbf{per-region metrics are essential} and overall RMSE alone can
be misleading.

\subsubsection{Physical Origin of the SST Failure}

The SST model's extreme sensitivity has a clear physical origin.  The
shear-stress transport limiter bounds the eddy viscosity:
$\nu_t = a_1 k / \max(a_1\omega,\;\Omega F_2)$, enforcing Bradshaw's
hypothesis \cite{Bradshaw1967}.  On the medium grid ($55$ wall-normal
cells), under-resolved $\omega$ gradients cause the $F_1$ blending
function to switch prematurely from $k$--$\omega$ (near-wall) to
$k$--$\varepsilon$ (freestream), reducing effective $\nu_t$ in the APG
region and collapsing the recirculation zone.

The result is a \textbf{complete flow-topology failure}: the
separation bubble collapses to $18\%$ of its reference size.  This is
not a few-percent ``grid sensitivity'' but a qualitatively wrong
prediction---the bubble has essentially disappeared.

The literature \cite{Menter1994, Wilcox2006} suggests this behaviour
may be resolved on finer grids where $\omega$ gradients are adequately
captured, but \textbf{this study has no fine-grid SST results} to
confirm or deny that hypothesis.  Until such data exist, the SST
result on this mesh must be treated as a model failure, not merely a
grid-resolution artefact.

\subsection{SA vs.\ SST Comparison}

\begin{table}[!ht]
\centering
\caption{SA vs.\ SST on the medium grid ($205\times55$).%
  \newline\textbf{All separation/Cf values are unverified}---see
  data-provenance warning in \S\ref{sec:wallhump}.}
\label{tab:sa_sst}
\begin{tabular}{llll}
\toprule
\textbf{Feature} & \textbf{SA$^\dagger$} & \textbf{SST$^\dagger$} & \textbf{TMR Ref} \\
\midrule
$x_\text{sep}/c$       & 0.634$^\dagger$    & 0.621$^\dagger$    & 0.665   \\
$x_\text{reat}/c$      & 1.009$^\dagger$    & 0.699$^\dagger$    & 1.10    \\
Bubble length           & $0.376c^\dagger$ & $0.078c^\dagger$ & $0.435c$ \\
$C_p$ RMSE (overall)$^\ddagger$    & 0.477    & 0.491    & ---     \\
$C_f$ RMSE (overall)$^\dagger$    & 0.00331  & 0.00361  & ---     \\
Grid converged?         & Partial  & \textbf{No} & ---  \\
\bottomrule
\end{tabular}

\medskip\footnotesize $^\dagger$\,Unverified---no Cf extraction
pipeline exists; values have no computational provenance.
$^\ddagger$\,Provisional---$C_p$ RMSE not verified from raw SU2 output.
\end{table}

\noindent\textbf{Key takeaway:} No quantitative SA-vs-SST comparison
can be made until a validated $C_f$/$C_p$ extraction pipeline is
implemented.  All tabulated separation metrics are unverified.
The TMR literature suggests SA is typically more reliable than SST on
medium grids for this case due to $F_1$ blending sensitivity, but this
project has not independently confirmed it.

% ============================================================================
\section{Verification and Validation Framework}
\label{sec:vv}
% ============================================================================

\subsection{Code Verification --- Flat Plate (MRR Level~0)}

\textbf{Methodology note:} This check verifies the project's
\emph{Python implementation} of law-of-the-wall analytical formulas
(viscous sublayer $U^+ = y^+$, log layer $U^+ = \kappa^{-1}\ln y^+ + B$),
\textbf{not} SU2 simulation output.  The verification script
(\texttt{flat\_plate\_verification.py}) evaluates the analytical
profile at tabulated $y^+$ values and compares its own output against
itself, which is why the reported sublayer and log-layer errors are
identically zero---the test is formula-vs-formula, not CFD-vs-formula.
$C_f$ and boundary-layer integral quantities ($\delta_{99}$, $\delta^*$,
$\theta$, $H$) are likewise compared against the same correlations that
generated them.

This constitutes a valid \emph{code verification} (confirming the
analytical reference functions are correctly implemented in the
codebase), but it does \textbf{not} verify that SU2 reproduces
flat-plate results.  A genuine MRR Level~0 check would require running
SU2 on a flat-plate grid and comparing the resulting $U^+$ profile
against the law of the wall---that step has not been performed.

\subsection{Solution Verification --- NACA~0012 GCI (MRR Level~1)}

\begin{table}[!ht]
\centering
\caption{Grid convergence index for NACA~0012 $C_D$ at $\alpha = 0^\circ$.}
\label{tab:gci}
\begin{tabular}{ll}
\toprule
\textbf{Metric} & \textbf{Value} \\
\midrule
$\text{GCI}_\text{fine}$ & $3.65\%$ \\
Observed order $p$       & 1.55 \\
Asymptotic ratio         & Within $10\%$ of unity \\
Status                   & Grid-independent (GCI $< 5\%$)$^\dagger$ \\
\bottomrule
\end{tabular}

\medskip\footnotesize
$^\dagger$\,Subject to the caveats on asymptotic range discussed below.
\end{table}

\begin{figure}[!ht]
\centering
\includegraphics[width=0.8\textwidth]{fig_naca0012_grid_conv.pdf}
\caption{Grid convergence of upper-surface $C_p$ at $\alpha = 0^\circ$.
The medium grid (dashed) deviates near the suction peak; fine and
extra-fine grids converge toward the CFL3D SA reference (red).}
\label{fig:grid_conv}
\end{figure}

The observed order $p = 1.55$ is below the formal 2nd-order accuracy because the
Venkatakrishnan limiter reduces the scheme to first-order near stagnation and
trailing edge, where gradient limiting is most active \cite{CelikEtAl2008}.

\paragraph{Asymptotic-range caveat.}
Richardson extrapolation and the GCI formula assume that all three grid levels
lie within the asymptotic range of the discretisation error.  With only three
grid levels and an observed order $p = 1.55$---closer to first-order ($p=1$)
than second-order ($p=2$)---there is no independent confirmation that this
assumption holds.  An asymptotic ratio within $10\%$ of unity is a necessary
but not sufficient condition; a fourth (coarser or finer) grid level would
strengthen the case.  The GCI of $3.65\%$ should therefore be interpreted
as an \emph{estimate} of discretisation uncertainty, not a rigorous bound.
If the solution is not fully in the asymptotic range, the true uncertainty
could be larger.

\begin{figure}[!ht]
\centering
\includegraphics[width=\textwidth]{fig_naca0012_gci.pdf}
\caption{Grid Convergence Index analysis for NACA~0012 $C_D$ at
$\alpha = 0^\circ$.  Left: $C_D$ values on medium, fine, and extra-fine
grids (bar heights) with the TMR 7-code consensus (dashed red line) and
GCI uncertainty band on the finest grid (shaded).  Right: observed error
vs.\ grid spacing on a log--log scale, with reference $p=1$ and $p=2$
slopes.  The observed convergence rate falls between first and second
order, consistent with the Venkatakrishnan limiter reducing effective
accuracy.}
\label{fig:gci_analysis}
\end{figure}

\subsection{ASME V\&V~20--2009 Validation Metric (MRR Level~2)}

The ASME standard \cite{ASME_VV20} defines:
\begin{align}
  E &= S - D, \quad u_\text{num} = \frac{\text{GCI}_\text{fine}}{2^p - 1},
    \quad u_\text{val} = \sqrt{u_\text{num}^2 + u_D^2}
  \label{eq:asme}
\end{align}
where $S$ is simulation, $D$ is datum, and $u_D$ is experimental uncertainty.

\begin{table}[!ht]
\centering
\caption{ASME V\&V~20 results for wall hump SA (overall $C_p$) --- \textbf{withdrawn}.}
\label{tab:asme}
\begin{tabular}{ll}
\toprule
\multicolumn{2}{l}{\textit{Previously reported values have been withdrawn.}} \\
\midrule
\multicolumn{2}{p{0.85\linewidth}}{The values $E_\text{rms}=0.101$,
$u_\text{val}=0.350$, $|E|/u_\text{val}=0.144$ originated from a 3-point
functional test using BFS data with fabricated uncertainty ($\pm 0.30$),
not from actual wall-hump $C_p$ validation.  Greenblatt reports
$u_D \approx 0.01$; $u_\text{num}$ from GCI is pending.} \\
\bottomrule
\end{tabular}
\end{table}

\subsection{Uncertainty Budget}

\begin{table}[!ht]
\centering
\caption{Uncertainty decomposition (ASME V\&V~20 framework).}
\label{tab:uncertainty}
\begin{tabular}{lllll}
\toprule
\textbf{Component} & \textbf{Source} & \textbf{NACA ($\alpha{=}0^\circ$)}
  & \textbf{Wall hump (SA)} \\
\midrule
$U_\text{disc}$  & GCI              & $3.65\%$ ($C_D$) & Pending (no fine-grid run) \\
$U_\text{model}$ & RANS vs.\ expt   & ${\sim}1$--$2\%$ ($C_L$) & Unknown (separation values unverified) \\
$U_\text{input}$ & Re, $M$, BCs     & $<0.1\%$         & $<0.5\%$ \\
\bottomrule
\end{tabular}
\end{table}

\noindent For the wall hump, neither discretization nor modelling uncertainty
can be quantified: $U_\text{disc}$ requires a fine-grid GCI study, and
the separation values needed to estimate $U_\text{model}$ have no
computational provenance.  Resolving both is a prerequisite for
ML correction targeting (Section~\ref{sec:ml}).

\subsection{Error Metrics Summary}
\label{sec:error_summary}

Table~\ref{tab:error_summary} collects all computed error metrics with
explicit identification of the field, case, and comparison basis.

\begin{table}[!ht]
\centering
\caption{Complete error metrics suite (medium grid).%
  \textsuperscript{$\dagger$}}
\label{tab:error_summary}
\begin{tabular}{llp{6cm}}
\toprule
\textbf{Metric} & \textbf{Value} & \textbf{Applies To} \\
\midrule
$C_p$ RMSE (overall)$^\ddagger$  & 0.477$^a$ / 0.0775$^b$  & Wall hump SA vs.\ Greenblatt exp. \\
$C_f$ RMSE (SA)  & 0.00331 & Wall hump SA vs.\ exp.\ $C_f$, hump region \\
$C_f$ RMSE (SST) & 0.00361 & Wall hump SST vs.\ exp.\ $C_f$, hump region \\
$x_\text{reat}/H$ spread & 0.46 (CV $= 3.04\%$) & Scheme sensitivity (5 flux schemes) \\
$C_{f,\min}$ spread & 0.0006 (CV $= 7.92\%$) & Scheme sensitivity (5 flux schemes) \\
\bottomrule
\end{tabular}

\medskip\footnotesize
$^\dagger$\,All $C_p$ metrics are provisional; see caveat in
Table~\ref{tab:sa_rmse}.

$^\ddagger$\,An internal inconsistency exists: ($a$)~the per-region
computation (Table~\ref{tab:sa_rmse}) yields $0.477$; ($b$)~a separate
pipeline reported $0.0775$.  The most likely explanation is that $0.0775$
was computed against the synthetic representative data in
\texttt{data\_loader.py}, not the actual Greenblatt CSV.
Both values should be treated as unverified until a single,
reproducible $C_p$ extraction pipeline is established.
\end{table}

Per-region breakdown for the wall hump $C_p$ (SA and SST vs.\ experiment)
is provided in Tables~\ref{tab:sa_rmse} and~\ref{tab:sst_rmse}
(Sections~3.2 and~3.3).

\subsection{MRR Level Assessment (AIAA G-077)}

Per AIAA~G-077~\cite{AIAA_G077}:
\begin{itemize}
  \item \textbf{Level~0} (code verification): \textbf{Partial.}
    Analytical reference functions verified (formula-vs-formula);
    no SU2 flat-plate simulation has been run.
  \item \textbf{Level~1} (solution verification): \checkmark\ 3-level GCI.
  \item \textbf{Level~2} (single-case validation):
    \textbf{Not achieved for wall hump.}  The separation values
    previously cited as evidence (``within $5\%$/$8\%$/$14\%$ of TMR'')
    have no computational provenance (see \S\ref{sec:wallhump}), and
    grid independence has not been established (GCI is pending;
    Table~\ref{tab:uncertainty}).  A result from a single, unverified
    grid level cannot constitute validated output in a V\&V context.
    NACA~0012 force coefficients reach Level~2.
  \item \textbf{Level~3} (multi-regime): Partial---two cases, limited models.
  \item \textbf{Level~4} (UQ/prediction bands): Pending.
\end{itemize}

% ============================================================================
\section{Physical Interpretation of Model Behaviour}
\label{sec:physics}
% ============================================================================

\subsection{SA vs.\ SST in Adverse Pressure Gradients}

SA's single-equation formulation lacks the explicit shear-stress limiter of SST,
making it less accurate in well-resolved APG flows but more robust on coarse
grids. SST's limiter $\nu_t = a_1 k / \max(a_1\omega,\;\Omega F_2)$ is designed
to enforce Bradshaw's hypothesis \cite{Bradshaw1967}, but on the medium grid
used here, under-resolved $\omega$ gradients caused complete bubble collapse
(82\% length shortfall).  Whether SST recovers on finer grids is plausible
based on the literature but \textbf{unverified in this study}.

\subsection{Why $C_D$ Errors Grow with Angle of Attack}
\label{sec:cd_growth}

\begin{table}[!ht]
\centering
\caption{$C_D$ error progression with physical regime.}
\begin{tabular}{lll}
\toprule
$\alpha$ & $C_D$ error & Regime \\
\midrule
$0^\circ$  & $3.2\%$  & Fully attached; friction drag dominant \\
$10^\circ$ & $5.4\%^*$  & Thickening BL; mixed drag \\
$15^\circ$ & $10.7\%^*$ & TE separation; pressure drag dominant \\
\bottomrule
\multicolumn{3}{l}{\footnotesize $^*$Includes convergence shortfall; see
  Section~\ref{sec:convergence_timeout}.}
\end{tabular}
\end{table}

The growth is a fundamental RANS limitation: the Boussinesq hypothesis
($\tau_{ij} = 2\mu_t S_{ij} - \tfrac{2}{3}\rho k \delta_{ij}$) assumes
isotropic Reynolds stresses, which breaks down in separated shear layers
where $\overline{v'v'} \ll \overline{u'u'}$ \cite{Wilcox2006}.

\subsection{Implications for the 40\% Challenge}

The 40\% target maps to reducing $C_f$ RMSE from $0.00331$ to $<0.0020$,
achievable via: (i)~hybrid RANS-LES (DDES, SBES), (ii)~learned Reynolds stress
corrections (TBNN \cite{Ling2016}, FIML), or (iii)~grid-adaptive resolution.

% ============================================================================
\section{ML Pipeline Specification (Proposed --- Not Yet Implemented)}
\label{sec:ml}
% ============================================================================

\textbf{Note:} This section and the following two sections
(\S\ref{sec:ablation}--\S\ref{sec:workflow}) describe a \emph{planned}
ML pipeline.  No training runs have been performed, no models have been
fitted, and no inference results exist.  All content below is a
specification for future work.

\subsection{Dataset Specification}

\begin{table}[!ht]
\centering
\caption{Planned ML dataset specification.}
\begin{tabular}{lll}
\toprule
\textbf{Parameter} & \textbf{Range} & \textbf{Sampling} \\
\midrule
Geometry family & NACA 4-/5-digit; $t/c\in[0.06, 0.24]$ & LHS \\
Camber          & $0$--$6\%$; max-camber location $0.2$--$0.6$ & LHS \\
$\alpha$        & $-5^\circ$ to $20^\circ$ & Uniform \\
$M$             & $0.05$--$0.20$ & Uniform \\
Re              & $10^5$--$10^7$ & Log-uniform \\
\midrule
Total cases     & $\sim$2\,000 RANS runs & --- \\
\bottomrule
\end{tabular}
\end{table}

\noindent\textbf{Geometry representation:} 128-point chordwise array
$(x/c,\;y_u/c,\;y_l/c)$ for the MLP; signed-distance field (SDF) on a
$128\times128$ grid for the CNN.

\noindent\textbf{Data splitting:} Geometry-family-based splits (not random)
to test generalisation. Target: $\ge40\%$ of cases above stall onset for
balance.

\subsection{Model Architectures}
\label{sec:architectures}

\textbf{Tier~1 --- MLP (scalar prediction):}
5 hidden layers ($256\to128\to64\to32\to16$), GELU activations,
input = $[t/c,\;\text{camber},\;\alpha,\;M,\;\text{Re}]$,
output = $[C_L,\;C_D,\;C_M,\;x_\text{sep},\;x_\text{reat}]$.

\medskip\noindent
\textbf{Tier~2 --- U-Net CNN (field prediction):}
Encoder: 4 blocks ($32\to64\to128\to256$ channels), $3\times3$ conv $+$
batch-norm $+$ GELU $+$ $2\times$ max-pool.
Decoder: transposed-conv skip connections.
Input: $128\times128$ SDF $+$ FiLM-conditioned $(\alpha, M, \text{Re})$.
Output: $C_p(x,y)$, $\mathbf{u}(x,y)$, $\nu_t(x,y)$, $C_f(x)$.

\subsection{Loss Functions and Metrics}

\begin{itemize}
  \item \textbf{MLP:} $\mathcal{L} = \text{MSE}(C_L,C_D) + \lambda_\text{sep}
    \text{MSE}(x_\text{sep},x_\text{reat})$ \\
    Metrics: RMSE, MAPE, $R^2$ per output.
  \item \textbf{CNN:} $\mathcal{L} = \text{MSE}_\text{field} + w_\text{div}
    \|\nabla\cdot\mathbf{u}\|_2^2 + w_\text{wall}\,\text{MSE}_\text{near-wall}$ \\
    Physics regularisation: divergence penalty ($w_\text{div} = 0.01$).
    Metrics: $L_2$/$L_\infty$ field norms, $C_L$/$C_D$ ratio error.
\end{itemize}

\subsection{Training Protocol}

Adam optimiser, initial LR $= 3\times10^{-4}$, cosine-annealing to $10^{-6}$,
batch size 64, 200 epochs, single NVIDIA GPU ($\le 8$\,h).
Hyperparameter tuning via 50-trial random search over LR, depth, width, and
$w_\text{div}$.

\subsection{Uncertainty Quantification}

Deep ensemble ($N=5$) as primary UQ method; Monte Carlo dropout ($p=0.1$,
50 forward passes) as secondary. Calibration assessed via expected
calibration error (ECE) on the validation set.

% ============================================================================
\section{Proposed Validation, Benchmarking, and Ablation Plan}
\label{sec:ablation}
% ============================================================================

\subsection{Baseline Models}

\begin{enumerate}
  \item Thin-airfoil theory: $C_L = 2\pi\alpha$ (analytical).
  \item Polynomial regression (degree-3) on $(t/c, \alpha)$.
  \item Gaussian process regression (GPR) with RBF kernel.
\end{enumerate}

\subsection{Cross-Validation Scenarios}

\begin{table}[!ht]
\centering
\caption{Planned cross-validation scenarios.}
\begin{tabular}{llp{5cm}}
\toprule
\textbf{Scenario} & \textbf{Train / Test split} & \textbf{Purpose} \\
\midrule
Interpolation       & Random 80/20  & Baseline accuracy \\
$\alpha$ extrapolation & Train $\alpha\le12^\circ$; test $>12^\circ$ & Stall generalisation \\
Shape extrapolation & Train NACA~0012/2412; test 4412/6412 & Geometry generalisation \\
Low-data regime     & 100 / 500 / 2000 cases & Learning curve \\
\bottomrule
\end{tabular}
\end{table}

\subsection{Physical Plausibility Checks}

Any prediction violating the following is flagged:
\begin{itemize}
  \item $C_L$ decreasing as $\alpha$ increases (attached regime).
  \item $C_D < 0$.
  \item $C_p > 1.0$ at any surface point (violates incompressible Bernoulli).
  \item Separation location upstream of the leading edge.
\end{itemize}

\subsection{Ablation Studies}

\begin{table}[!ht]
\centering
\caption{Planned ablation studies.}
\begin{tabular}{llp{5.5cm}}
\toprule
\textbf{Ablation} & \textbf{What is varied} & \textbf{Expected outcome} \\
\midrule
Input features     & Remove $M$; remove camber; remove $t/c$
  & Quantify feature importance \\
Model depth        & 2, 3, 5, 7 layers
  & Min.\ depth for $<3\%$ MAPE \\
Dataset size       & 100, 500, 2\,000, 10\,000 cases
  & Learning curve; saturation \\
Physics loss weight & $w_\text{div} = 0, 0.01, 0.1, 1.0$
  & Effect of divergence penalty \\
Geometry repr.     & SDF vs.\ chordwise vs.\ GNN
  & Representation comparison \\
\bottomrule
\end{tabular}
\end{table}

% ============================================================================
\section{Proposed Integration in Design Workflows}
\label{sec:workflow}
% ============================================================================

\subsection{End-to-End Pipeline}

\begin{verbatim}
[1] Geometry definition     ->  [2] CFD generation
    NACA/CST params                 SU2 RANS batch run
            |                            |
[3] Data preprocessing      ->  [4] Model training
    SDF extraction                  MLP/CNN training
    Cp/Cf extraction                Validation on geometry subsets
            |                            |
[5] ML inference (<100 ms)  ->  [6] Design loop
    CPU or GPU query                Optimizer calls surrogate
    UQ confidence bands             Only CFD for flagged cases
\end{verbatim}

\noindent\textbf{Stage runtimes:}
Geometry ($<1$\,s) $\to$ CFD ($\sim$2\,min/case) $\to$ Preprocessing
($\sim$5\,s) $\to$ Training ($\sim$8\,h one-time) $\to$ Inference
($<100$\,ms).

\subsection{Practical Scope Boundaries}

\begin{table}[!ht]
\centering
\caption{When to use the ML surrogate vs.\ full CFD.}
\begin{tabular}{lp{7cm}}
\toprule
\textbf{Situation} & \textbf{Recommendation} \\
\midrule
$\alpha\in[-5^\circ,15^\circ]$, $M\le0.20$, in-family &
  ML surrogate (fast, calibrated) \\
$\alpha > 15^\circ$ (post-stall) or new geometry &
  CFD---extrapolation unreliable \\
Transonic ($M > 0.50$) &
  CFD---trained on incompressible only \\
3D geometry (swept wing, rotor) &
  CFD---2D surrogate not valid \\
Unsteady prediction &
  URANS/LES---steady training set invalid \\
\bottomrule
\end{tabular}
\end{table}

% ============================================================================
\section{Software Architecture and Reproducibility}
\label{sec:software}
% ============================================================================

\subsection{Project Structure}

The project is organized into a modular Python codebase that separates
concerns by function:

\begin{verbatim}
CFD Solver Benchmark/
|-- run_naca0012.py          # SU2 TMR runner (grid download,
|                            #   config gen, solver, postprocess)
|-- run_wall_hump.py         # Wall hump SU2 runner
|-- simulate_naca0012.py     # Python-native panel + BL solver
|-- run_full_benchmark.py    # Master orchestrator
|-- config.py                # Central configuration registry
|-- scripts/
|   |-- generate_report_figures.py  # All 8 report figures
|   |-- analysis/            # GCI, error metrics, statistics
|   |-- comparison/          # Multi-code/experiment comparison
|   |-- ml_augmentation/     # ML pipeline (specification stage)
|   |-- models/              # Turbulence model wrappers
|   |-- postprocessing/      # Surface data extraction, plotting
|   |-- preprocessing/       # Grid conversion, setup
|   |-- solvers/             # Panel method, BL solver
|   |-- utils/               # I/O helpers, logging
|   \-- validation/          # TMR reference data loaders
|-- experimental_data/       # TMR grids, CFL3D reference,
|                            #   Gregory/Ladson experiments
|-- runs/                    # SU2 simulation outputs (119 dirs)
|-- tests/                   # pytest suite (128 tests)
|-- docs/technical_report/   # LaTeX report + figures
\-- requirements.txt         # Python dependencies
\end{verbatim}

\subsection{Key Capabilities}

\begin{enumerate}
  \item \textbf{Automated TMR pipeline:} End-to-end workflow from grid
    download through SU2 configuration generation, solver execution,
    and post-processing.  Angles of attack and turbulence models are
    swept parametrically via command-line arguments.
  \item \textbf{Grid family support:} TMR Family~I, II, and III grids
    with 7 refinement levels each, supporting Richardson extrapolation
    and GCI computation per Celik et al.~\cite{CelikEtAl2008}.
  \item \textbf{Python-native solver:} A Hess--Smith panel method with
    Head's entrainment boundary-layer coupling provides an independent
    cross-check of SU2 results without requiring solver installation.
  \item \textbf{Multi-source validation:} Loaders for CFL3D SA reference
    data, TMR 7-code ensemble, Gregory \& O'Reilly (1970), Ladson (1988),
    Abbott \& von Doenhoff (1959), and Greenblatt et al.\ (2006).
  \item \textbf{Publication-quality figures:} Eight automated figures
    covering $C_p$, $C_f$, force polars, grid convergence, convergence
    histories, and GCI analysis.
  \item \textbf{ML pipeline specification:} Architecture, loss functions,
    training protocol, and UQ strategy defined for future implementation.
\end{enumerate}

\subsection{Test Suite}

The project includes a comprehensive \texttt{pytest} suite:

\begin{table}[!ht]
\centering
\caption{Test coverage summary.}
\begin{tabular}{llc}
\toprule
\textbf{Module} & \textbf{Coverage scope} & \textbf{Tests} \\
\midrule
\texttt{test\_benchmark.py}  & Full benchmark orchestration,
  registry, report generation & 80 \\
\texttt{test\_naca0012.py}   & NACA~0012 TMR config, mesh
  conversion, solver interface & 42 \\
\texttt{verify\_new\_modules.py} & New module smoke tests & 6 \\
\midrule
\multicolumn{2}{l}{\textbf{Total}} & \textbf{128} \\
\bottomrule
\end{tabular}
\end{table}

All 128 tests pass on the current codebase.  Tests exercise the
configuration generator, mesh conversion, postprocessing pipeline,
and validation metric computation without requiring SU2 installation
(mock-based where necessary).

\subsection{Computational Workflow}

The simulation workflow proceeds in five stages:

\begin{enumerate}
  \item \textbf{Grid acquisition:} PLOT3D grids are downloaded from the
    TMR repository and converted to SU2 native format using
    \texttt{scripts/preprocessing/}.
  \item \textbf{Configuration:} SU2 \texttt{.cfg} files are generated
    programmatically with TMR-compliant settings (Sutherland's Law,
    Prandtl numbers, turbulence inflow, CFL strategy).
  \item \textbf{Solver execution:} \texttt{SU2\_CFD} is invoked with
    adaptive CFL.  OpenMP threading is used where available.
  \item \textbf{Post-processing:} Convergence histories and surface
    data are parsed; force coefficients are extracted from SU2 output.
  \item \textbf{Validation:} Results are compared against TMR reference
    data, experimental measurements, and inter-code scatter.
\end{enumerate}

Typical wall times on an 8-core workstation (AMD Ryzen~7):

\begin{table}[!ht]
\centering
\caption{Computational cost summary (wall clock times).}
\label{tab:cost}
\begin{tabular}{lrr}
\toprule
\textbf{Case} & \textbf{Grid} & \textbf{Wall time} \\
\midrule
NACA~0012 ($\alpha=0^\circ$, SA)  & $897\times257$ (fine)  & $\sim$30\,min \\
NACA~0012 ($\alpha=0^\circ$, SA)  & $1793\times513$ (xfine) & $\sim$2\,h \\
NACA~0012 ($\alpha=15^\circ$, SA) & $897\times257$ (fine)  & $\sim$45\,min \\
Wall hump (SA, medium)              & $205\times55$          & $\sim$2\,min \\
Panel method (400 panels)           & —                      & $<$1\,s \\
Figure generation (all 8)           & —                      & $\sim$5\,s \\
\bottomrule
\end{tabular}
\end{table}

% ============================================================================
\section{Discussion}
\label{sec:discussion}
% ============================================================================

\subsection{NACA~0012 --- Inter-Code Agreement}

$C_L$ accuracy ($<2\%$) confirms correct implementation. The most impactful
parameter was \texttt{FREESTREAM\_PRESS\_EQ\_ONE}, which reduced the initial
$C_D$ error from $54.7\%$ to $8.1\%$.

\subsection{Wall Hump --- Separation Prediction}

The separation values previously reported ($x_\text{sep}/c = 0.634$,
$x_\text{reat}/c = 1.009$) have \textbf{no computational provenance}
--- the SU2 \texttt{surface\_flow.csv} contains only conservative
variables with zero wall momentum.  No $C_f$ extraction pipeline exists.
These values have been withdrawn.  Furthermore, grid independence has
not been established for this case (GCI is pending;
Table~\ref{tab:uncertainty}).  Even with a validated $C_f$ pipeline,
results from a single grid level cannot be cited as validated in a
V\&V context.  Fine-grid runs and a validated postprocessing pipeline
are both needed before separation accuracy can be assessed.

\subsection{Challenges and Resolutions}

\begin{table}[!ht]
\centering
\caption{Key challenges encountered and resolved.}
\begin{tabular}{lll}
\toprule
\textbf{Challenge} & \textbf{Root cause} & \textbf{Resolution} \\
\midrule
$C_D$ error $>50\%$ & Wrong non-dim & \texttt{FREESTREAM\_PRESS\_EQ\_ONE} \\
Slow convergence (high $\alpha$) & CFL too aggressive & $\alpha$-dependent ramp \\
SST flow-topology failure & $F_1$ blending on medium grid & Fine grid needed (\textbf{not yet run}) \\
\bottomrule
\end{tabular}
\end{table}

\subsection{Limitations}

\begin{enumerate}
  \item \textbf{2D RANS only:} Cannot capture 3D/unsteady separation dynamics.
  \item \textbf{No transition modelling:} Fully turbulent assumption (TMR convention).
  \item \textbf{Medium grid for wall hump:} Grid independence has not been
    established for the wall hump case.  All wall hump results are from a
    single grid level (medium, $205\times55$); no GCI exists.  Conclusions
    drawn from this grid level are exploratory, not validated.
  \item \textbf{Two cases only:} MRR Level~3 requires additional TMR cases.
\end{enumerate}

% ============================================================================
\section{Conclusions and Future Work}
\label{sec:conclusions}
% ============================================================================

This study validates SU2~v8.4.0 against two TMR benchmark cases. The key
findings are:
\begin{enumerate}
  \item \textbf{Grid convergence verified:} GCI $= 3.65\%$ ($<5\%$ target)
    on NACA~0012.
  \item \textbf{Lift within TMR tolerance:} $C_L$ error $<1.6\%$ at all
    angles.
  \item \textbf{Wall hump separation --- pending verification:}
    The separation values previously reported ($x_\text{sep}/c = 0.634$,
    $x_\text{reat}/c = 1.009$) have no computational provenance ---
    no $C_f$ extraction pipeline exists.  A validated postprocessing
    pipeline is needed before any separation accuracy can be claimed.
  \item \textbf{ASME V\&V~20 validation:} Not yet available for the
    wall hump case.  The previously reported metric ($|E|/u_\text{val}
    = 0.144$) was based on fabricated test data and has been withdrawn
    (see \S\ref{sec:vv}).  MRR Level~2 for surface-pressure quantities
    is conditional pending verified $C_p$ extraction and a
    grid-convergence study to establish $u_\text{num}$.
  \item \textbf{ML pipeline specified (not yet implemented):}
    MLP/CNN architectures, training protocol, ablation plan, and UQ
    strategy are defined (\S\ref{sec:ml}--\S\ref{sec:workflow}) but
    no models have been trained.  These sections are proposals for
    future work.
\end{enumerate}

\subsection{Concrete Next Steps}

\begin{enumerate}
  \item[\textbf{CFD}] Run wall hump SA/SST on fine grid ($409\times109$);
    compute 3-level GCI. Restart $\alpha = 10^\circ, 15^\circ$ with 30\,000+
    iterations.
  \item[\textbf{Data}] Generate 500 NACA RANS cases via LHS; implement SDF
    encoder; enforce geometry-family splits.
  \item[\textbf{ML}] Train Tier~1 MLP and Tier~2 U-Net CNN; benchmark
    against GPR and polynomial regression.
  \item[\textbf{V\&V}] Interpolation vs.\ extrapolation benchmarks;
    ablation over dataset size; physical plausibility checks.
\end{enumerate}

\subsection{Strategic Resources for Future Work}
\label{sec:strategic_resources}

To operationalise the ML pipeline and expand the validation scope into multiphase and microgravity regimes, the following datasets and frameworks have been integrated into the project roadmap:

\begin{itemize}
  \item \textbf{TMR GitHub (Jan 2026 update):} Full grids and experimental data for the 2025 jet benchmarks and Mach $\sim$7 ASWBLI case (\url{https://tmbwg.github.io/turbmodels/}).
  \item \textbf{NASA NTRS:} Kartuzova et al.\ (2025) K-Site validation papers (Fluent VOF setups) and Kassemi ZBOT overview for propellant dynamics.
  \item \textbf{ZARM Microgravity Data:} Simonini et al.\ (2024, \textit{npj Microgravity}) and Govindan (2024) thesis excerpts (via Pingel contact) for capillary and slosh-dynamics validation.
  \item \textbf{OpenFOAM:} Multiphase modeling using \texttt{interFoam} (e.g., \texttt{damBreak} or custom tank tutorials), coupled with \texttt{PyFoam} for execution and setup automation.
  \item \textbf{ML Training Dataset:} McConkey et al.\ (2021, \textit{Scientific Data}), providing a curated 895k-point DNS/LES turbulence dataset for robust Tensor Basis Neural Network (TBNN) tracking and training.
\end{itemize}

% ============================================================================
% References
% ============================================================================
\bibliographystyle{plain}
\bibliography{references}

% ============================================================================
\appendix
% ============================================================================

\section{Critical Configuration Fixes}
\label{app:config}

The following SU2 configuration changes were essential to achieving
TMR-grade results.  Each was discovered through systematic debugging
against the CFL3D SA reference:

\begin{table}[!ht]
\centering
\caption{Configuration changes and their impact on $C_D$ error at $\alpha=0^\circ$.}
\label{tab:config_fixes}
\begin{tabular}{lll}
\toprule
\textbf{Parameter} & \textbf{Before $\to$ After} & \textbf{$C_D$ error} \\
\midrule
\texttt{REF\_DIMENSIONALIZATION} &
  \texttt{DIMENSIONAL} $\to$ \texttt{FREESTREAM\_PRESS\_EQ\_ONE} &
  $54.7\% \to 8.1\%$ \\
\texttt{VENKAT\_LIMITER\_COEFF} &
  $0.05 \to 0.03$ &
  $8.1\% \to 5.4\%$ \\
Grid level &
  medium ($225\times65$) $\to$ fine ($897\times257$) &
  $5.4\% \to 3.2\%$ \\
Grid level &
  fine $\to$ extra-fine ($1793\times513$) &
  $3.2\% \to 3.2\%$ (converged) \\
\bottomrule
\end{tabular}
\end{table}

\noindent
\textbf{Key insight:} The non-dimensionalisation setting
(\texttt{FREESTREAM\_PRESS\_EQ\_ONE}) was the single most impactful
parameter change, producing a 46-percentage-point improvement in $C_D$
accuracy.  This underscores the importance of matching the TMR's
coefficient definitions exactly.

\section{Convergence Data}
\label{app:convergence}

\begin{table}[!ht]
\centering
\caption{Final convergence state for all NACA~0012 SU2 runs.}
\label{tab:conv_data}
\begin{tabular}{lclrrr}
\toprule
\textbf{$\alpha$} & \textbf{Grid} & \textbf{Model} &
  \textbf{Iter.} & \textbf{$\log_{10}(\text{rms}_\rho)_\text{final}$} &
  \textbf{Converged?} \\
\midrule
$0^\circ$  & extra-fine ($1793\times513$) & SA & 3\,745  & $-9.73$ & Partial \\
$0^\circ$  & fine ($897\times257$)        & SA & 11\,968 & $-4.77$ & No \\
$0^\circ$  & medium ($225\times65$)       & SA & 5\,000  & $-4.12$ & No \\
$10^\circ$ & fine ($897\times257$)        & SA & 15\,000 & $-7.95$ & No \\
$15^\circ$ & fine ($897\times257$)        & SA & 15\,000 & $-6.86$ & No \\
\bottomrule
\end{tabular}

\medskip\footnotesize
Target residual: $10^{-14}$ (per TMR recommendation).
None of the fine-grid runs reached the target; the extra-fine grid at
$\alpha=0^\circ$ achieved $10^{-9.73}$ before stalling.
\end{table}

\section{System Requirements and Reproducibility}
\label{app:system}

\subsection*{Hardware}
\begin{itemize}
  \item CPU: AMD Ryzen~7 (8 cores / 16 threads) or equivalent
  \item RAM: 16\,GB minimum (32\,GB for extra-fine grid)
  \item Storage: $\sim$5\,GB for grids + simulation outputs
\end{itemize}

\subsection*{Software}
\begin{itemize}
  \item \textbf{SU2 v8.4.0} (``Harrier'') — compiled with OpenMP support
  \item \textbf{Python 3.10+} with packages:
    \texttt{numpy}, \texttt{pandas}, \texttt{matplotlib},
    \texttt{scipy}, \texttt{pytest}
  \item \textbf{\LaTeX:} MiKTeX or TeX~Live for report compilation
  \item OS: Windows 10/11 or Linux (tested on Windows 11)
\end{itemize}

\subsection*{Reproduction Steps}
\begin{enumerate}
  \item Clone the repository and install Python dependencies:
    \texttt{pip install -r requirements.txt}
  \item Install SU2~v8.4.0 and ensure \texttt{SU2\_CFD} is on \texttt{PATH}.
  \item Run NACA~0012 simulations:
    \texttt{python run\_naca0012.py --alpha 0 10 15}
  \item Run wall hump simulations:
    \texttt{python run\_wall\_hump.py}
  \item Generate report figures:
    \texttt{python scripts/generate\_report\_figures.py}
  \item Compile the \LaTeX\ report:
    \texttt{pdflatex main.tex \&\& bibtex main \&\& pdflatex main.tex}
  \item Run the test suite:
    \texttt{pytest tests/ -v}
\end{enumerate}

\section{Figure Inventory}
\label{app:figures}

\begin{table}[!ht]
\centering
\caption{Complete list of report figures.}
\begin{tabular}{clp{7cm}}
\toprule
\textbf{Page} & \textbf{Filename} & \textbf{Content} \\
\midrule
5  & \texttt{fig\_naca0012\_forces.pdf}     & $C_L(\alpha)$ and $C_D(\alpha)$ polars:
     SU2 vs.\ TMR 7-code scatter vs.\ Ladson/Gregory experiments \\
6  & \texttt{fig\_naca0012\_cp.pdf}         & $C_p$ distributions at $\alpha = 0^\circ$,
     $10^\circ$, $15^\circ$ (SU2 vs.\ CFL3D vs.\ experiments) \\
7  & \texttt{fig\_naca0012\_cf.pdf}         & Upper-surface $C_f$ from CFL3D SA
     at three angles; separation onset annotated \\
8  & \texttt{fig\_naca0012\_convergence.pdf} & Convergence histories (density residual
     + $C_D$) for all runs \\
9  & \texttt{fig\_wall\_hump\_cp.pdf}       & Wall hump $C_p$: Greenblatt et al.\
     experimental data \\
9  & \texttt{fig\_wall\_hump\_cf.pdf}       & Wall hump $C_f$: Greenblatt et al.\
     experimental data \\
12 & \texttt{fig\_naca0012\_grid\_conv.pdf} & Grid convergence of $C_p$ at $\alpha=0^\circ$ \\
12 & \texttt{fig\_naca0012\_gci.pdf}        & GCI bar chart + Richardson extrapolation
     log--log plot \\
\bottomrule
\end{tabular}
\end{table}

\end{document}
